\documentclass[11pt]{article}
\usepackage{fancyhdr}
\usepackage[utf8]{inputenc}	% Para caracteres en español
\usepackage{amsmath,amsthm,amsfonts,amssymb,amscd}
\usepackage{multirow,booktabs}
\usepackage[table]{xcolor}
\usepackage{fullpage}
\usepackage{lastpage}
% \usepackage{enumitem}
\usepackage{fancyhdr}
\usepackage{mathrsfs}
\usepackage{wrapfig}
\usepackage{setspace}
\usepackage{calc}
\usepackage{multicol}
\usepackage{cancel}
\usepackage[retainorgcmds]{IEEEtrantools}
\usepackage[margin=3cm]{geometry}
\usepackage{amsmath}
\usepackage{tikz}
\usepackage{blindtext}
\newlength{\tabcont}
\setlength{\parindent}{0.0in}
% \setlength{\parskip}{0.05in}
\usepackage{empheq}
\usepackage{framed}
\usepackage[most]{tcolorbox}
\usepackage{xcolor}
\usepackage{enumerate}
\colorlet{shadecolor}{gray!15}
\parindent 0in
% \parskip 12pt
\geometry{margin=1in, headsep=0.25in}

\theoremstyle{definition}
\newtheorem{defn}{Definition}
\newtheorem{reg}{Rule}
\newtheorem{exer}{Exercise}
\newtheorem{note}{Note}
\newenvironment{prf}[1][Proof]{\textbf{#1.} }{\qed}
\newtheorem{theorem}{Theorem}[section]
\newtheorem{corollary}{Corollary}[section]

\lhead{\includegraphics[height=1.0cm]{../GLE.pdf} GLE}
\chead{Cauchy-Schwarz Master Class}
\rhead{Encontro - 1}

\title{Chapter 1 - Review}
\author{Davi Barreira}
\begin{document}
\maketitle
\thispagestyle{fancy}
\setcounter{section}{0}
\section{Starting with Cauchy}

In these notes, we consider all
$a_i, b_j$ for $i,j \in \mathbb N$ to be real values.
For a vector space $V$ over $\mathbb R$, the inner product
$\langle \cdot, \cdot \rangle: V \times V \to \mathbb R$ has the following
properties:
\begin{enumerate}[(i)]
    \item $\langle \mathbf v, \mathbf v \rangle \geq 0$ and equal to $0$ iff $\mathbf v = 0$;
    \item $\langle \alpha \mathbf v + \mathbf w, \mathbf v \rangle
    = \alpha \langle \mathbf v, \mathbf v \rangle + \langle \mathbf w, \mathbf v \rangle$;
    \item $\langle \mathbf v, \mathbf w \rangle = \langle \mathbf w, \mathbf v \rangle$.
\end{enumerate}

\begin{theorem}{(Cauchy-Schwarz Inequality \textit{Finite} terms)}
    
   \begin{equation}
    (a_1 b_1 + \dots + a_n b_n)^2 \leq
    (a_1^2 +  \dots + a_n^2)(b_1^2  + \dots + b_n^2).
   \end{equation}
\end{theorem}
\begin{shaded}
\begin{prf}
    Using induction.
    \begin{equation*}
        (a_1 - b_1)^2 \geq 0 \implies a_1^2 + b_1^2 \geq 2 a_1 b_1.
    \end{equation*}
    Suppose that for $n$:
    \begin{equation*}
        \sum^n_{i=1} a_i b_i \leq
        \sqrt{\sum^n_{i=1}a_i^2}
        \sqrt{\sum^n_{i=1}b_i^2}.
    \end{equation*}
    \begin{equation*}
        \sum^n_{i=1} a_i b_i
        + a_{n+1} b_{n+1}
        \leq
        \sqrt{\sum^n_{i=1}a_i^2}
        \sqrt{\sum^n_{i=1}b_i^2}
        + a_{n+1} b_{n+1}.
    \end{equation*}

    Apply Cauchy-Schwarz to $c_1 := \sqrt{\sum^n_{i=1}a_i^2}$, $d_1 := \sqrt{\sum^n_{i=1}b_i^2}$,
    $c_2 := a_{n+1}$ and $d_2 := b_{n+1}$, hence:
    \begin{equation*}
        c_1 d_1 + c_2 d_2 \leq
        \sqrt{c_1 ^2 + c_2^2}
        \sqrt{d_1 ^2 + d_2^2} =
        \sqrt{\sum^n_{i=1}a_i^2 + a_{n+1}^2}
        \sqrt{\sum^n_{i=1}b_i^2 + b_{n+1}^2}.
    \end{equation*}
\end{prf}
\end{shaded}

\newpage
\begin{theorem}{(Cauchy-Schwarz Inequality \textit{Infinite} terms)}
   \begin{equation}
    \left( 
        \sum^{\infty}_{i=1} a_i b_i
    \right) ^2
    \leq
    \left( 
        \sum^{\infty}_{i=1}a_i^2
    \right)
    \left( 
    \sum^{\infty}_{i=1}b_i^2
    \right).
   \end{equation}
\end{theorem}

\begin{shaded}
\begin{prf}
    First, note that $(a+b)^2 \geq 0$, $(a-b)^2 \geq 0$, then
    $a^2 + b^2 \geq 2 |ab|$. Hence,
    \begin{equation*}
        \lim_{n\to \infty}
        \sum^n_{i=1} a_i ^2 +
        \sum^n_{i=1} b_i ^2
         = 
        \sum^{\infty}_{i=1} a_i ^2  
        +
        \sum^{\infty}_{i=1} b_i ^2
        \geq
        \lim_{n\to \infty}
        \sum^n_{i=1}2|a_i b_i| =
        \sum^{\infty}_{i=1} 2 |a_ib_i|.
    \end{equation*}
    Since the inequality above is valid for any $a_i,b_i \in \mathbb R$,
    make:
    \begin{equation*}
        \hat a_i := \frac{a_i}{\sqrt{\sum^{\infty}a_i^2}} ,\quad
        \hat b_i := \frac{b_i}{\sqrt{\sum^{\infty}b_i^2}}.
    \end{equation*}
    Note that
    \begin{equation*}
    \sum^{\infty}_{i=1} \hat a_i ^2 = 
    \frac{\sum^{\infty}_{i=1} a_i ^2}{\sum^{\infty}_{i=1} a_i ^2} = 1
    , \quad
    \sum^{\infty}_{i=1} \hat b_i ^2 = 
    \frac{\sum^{\infty}_{i=1} b_i ^2}{\sum^{\infty}_{i=1} b_i ^2} = 1.
    \end{equation*}
    Therefore,
    \begin{align*}
        2 =
        \sum^{\infty}_{i=1} \hat a_i ^2  
        +
        \sum^{\infty}_{i=1} \hat b_i ^2
        \geq
        \sum^{\infty}_{i=1} 2 |\hat a_i \hat b_i|
        \implies
        1 \geq 
        \sum^{\infty}_{i=1} |\hat a_i \hat b_i|
        = 
        \frac{\sum^{\infty}_{i=1} a_i b_i}
        {\sum^{\infty}_{i=1}a_i^2
        \sum^{\infty}_{i=1}b_i^2}.
    \end{align*}
\end{prf}
\end{shaded}

\begin{corollary}{(Sharpness of Cauchy-Schwarz Inequality}
    The Cauchy-Schwarz inequality is sharp iff
    $a_k = \lambda b_k$ for every $k \in \mathbb N$, where
    \begin{equation}
        \lambda = \frac{\sqrt{\sum^\infty_{i=1}a_i ^2}}
        {\sqrt{\sum^\infty_{i=1}b_i ^2}}.
    \end{equation}
\end{corollary}
\begin{shaded}
\begin{prf}
    Using the previous demonstration, we can see that
    \begin{align*}
        (\hat a_k - \hat b_k)^2 = 0 &\iff \hat a_k = \hat b_k \iff
        \frac{a_k}{\sqrt{\sum^\infty_{i=1} a_i^2}}=
        \frac{b_k}{\sqrt{\sum^\infty_{i=1} b_i^2}}\\
        &\iff
        a_k = \frac{\sqrt{\sum^\infty_{i=1}a_i^2}}
        {\sqrt{\sum^\infty_{i=1}b_i^2}} b_k.
    \end{align*}
\end{prf}
\end{shaded}
\end{document}